\documentclass[10pt]{article}
\usepackage[utf8]{inputenc}
\usepackage[T1]{fontenc}
\usepackage{amsmath}
\usepackage{amsfonts}
\usepackage{amssymb}
\usepackage[version=4]{mhchem}
\usepackage{stmaryrd}
\usepackage{bbold}

\begin{document}
LIG

L16: The order of an integer continued

\begin{itemize}
  \item We start this lecture with a sequence of results that establishes the number of primitive roots modulo $m$ (under the assumption that at least one primitive root modulo $m$ exists).\\
Proposition 5.4: Let $a, m \in \mathbb{Z}$ with $m>0$ and $(a, m)=1$. If $i$ is a positive integer, then
\end{itemize}

$$
\operatorname{ord}_{m}\left(a^{i}\right)=\frac{\operatorname{ord}_{m}(a)}{\left(\operatorname{ord}_{m}(a), i\right)} .
$$

Proof: Put $d=\left(\operatorname{or} d_{m} a, i\right)$. Then there exist $b, c \in \mathbb{Z}$ such that ord $m a=d b, i=d c$, and $(b, c)=1$. Now,\\
$\left(a^{i}\right)^{b}=\left(a^{d c}\right)^{(\operatorname{erd} m a) / d}=\left(a^{c}\right)^{\operatorname{ordm} a}=\left(a^{\operatorname{ord} m}\right)^{c}$

$$
\equiv 1(\bmod m) .
$$

By Proposition 5.1, we have that $\operatorname{ord}_{m}\left(a^{i}\right) \mid b$. Also, $a^{i\left(\operatorname{ord}_{m}\left(a^{i}\right)\right)}=\left(a^{i}\right)^{\operatorname{ord} m\left(a^{i}\right)} \equiv 1(\bmod m)$.

By Proposition 5.1, we have that $\operatorname{ord}_{m} a \mid i\left(\operatorname{ord} d_{m}\left(a^{i}\right)\right)$. Thus $d b \mid d c\left(\operatorname{ord} d_{m}\left(a^{i}\right)\right)$, and we have that $b \mid \operatorname{cord} d_{m}\left(a^{i}\right)$. Since $(b, c)=1$, we obtain blordm $\left(a^{i}\right)$. Thus

$$
\operatorname{ord}_{m}\left(a^{i}\right)=b=\frac{\operatorname{ord}_{m} a}{d}=\frac{\operatorname{ord}_{m} a}{(\operatorname{ord} a, i)},
$$

as desined.\\
Corollary 5.5: Let $a, m \in \mathbb{Z}$ with $m>0$ and $(a, m)=1$. If $i$ is a positive integer, then $\operatorname{ord}_{m}\left(a^{i}\right)=\operatorname{ord}_{m}(a)$ if and only if $(\operatorname{ord} m(a), i)=1$.\\
Proof: Apply Proposition 5.4.\\
Corollary 5.6: If a primitive root modulo $m$ exists, then there are exactly $\phi(\phi(m))$ incongruent primitive roots modulo $m$. Proof: Let $r$ be a primitive root modulo $m$.

Then $\operatorname{ord}_{m}{ }^{r}=\phi(m)$. By Proposition 5.3, (3) each residue class coprime to $m$ occurs exactly once in the set $\left\{r, r^{2}, r^{3}, \ldots, r^{\phi(m)}\right\}$. If $i$ is a positive integer such that $1 \leq i \leq \phi(m)$, then $\operatorname{ord} \boldsymbol{d}_{m}\left(r^{i}\right)=\operatorname{ord} \boldsymbol{d}_{m}(r)=\phi(m)$ if and only if $(\phi(m), i)=1$ by Corollary 5.5. The number of such integers $i$ with $1 \leqslant i \leqslant \phi(m)$ is $\phi(\phi(m))$, as desired.\\
Example: - We showed in a previous lecture that 3 is a primitive root modulo 7 .\\
Thus there are exactly $\phi(\phi(7))=\phi(6)=2$ incongruent primitive roots modulo 7 by Corollary 5.6. By Corollary 5.5,\\
ord, $\left(3^{i}\right)=$ ord, (3) if and only if $(\operatorname{ord}, 3, i)=1$.\\
Since ord, $(3)=6$, we must have $i=1,5$. Since\\
$3^{5} \equiv 5(\bmod 7), 5$ is the other incongrnent\\
primitive root modulo 7 .

\begin{itemize}
  \item In a previous lecture we proved that 2 was a primitive root modulo 13 . Thus there are exactly $\phi(\phi(13))=\phi(12)=4$ incongruent primitive roots modulo 13 by Corollary S.6.\\
Exercise: Find the other primitive roots.\\
Remark: Consider the following question:\\
"Since $\phi(\phi(8))=\phi(4)=2$, do there exist two primitive roots modulo 8 by corollary 5.6? The answer is no! We proved last lecture that there are no primitive roots modulo 8 . Corollary 5.6 does not apply here!
\end{itemize}

\section*{Primitive Roots for Prime Numbers}
\begin{itemize}
  \item Our next goal is to prove that primitive roots exist for prime numbers. We will need Lagrange's Theorem for this. Lagrange's Theorem concerns the number of solutions a polynomial can have modulo $P$.
\end{itemize}

Theorem 5.7 (Lagrange): Let $p$ be a prime 5 number and let

$$
f(x)=a_{n} x^{n}+a_{n-1} x^{n-1}+\ldots+a_{1} x+a_{0}
$$

be a polynomial of degree $n \geqslant 1$ with integral coefficients such that $p+a_{n}$. Then the congruence

$$
f(x) \equiv 0(\bmod p)
$$

has at most $n$ incongruent solutions modulo $p$.\\
Proof: We use induction on $n$. If $n=1$, we have $f(x)=a_{1} x+a_{0}$ with $p+a_{1}$. Since $\left(p, a_{1}\right)=1$, the linear congruence $a, x \equiv-a_{0}(\bmod p)$ has exactly one incongruent solution modulo $p$ by Theorem 2.6 . Thus $f(x) \equiv 0$ (modp) has exactly one incongruent solution modulo $p$, and the Theorem holds for $n=1$.\\
Assume that $k \geqslant 1$ and that the Theorem holds for polynomials of degree $n=k$. We must show that the Theorem holds for polynomials of degree\\
$n=k+1$. Accordingly, let

$$
f(x)=a_{k+1} x^{k+1}+a_{k} x^{k}+\cdots+a_{1} x+a_{0}
$$

be a polynomial of degree $k+1$ with integral coefficients such that $p+a_{k H}$. If $f(x) \equiv 0(\bmod p)$ has no solutions, then the Theorem holds and the proof is complete. Assume the congruence has at least one solution, say a. View $f(x) \in \mathbb{F}_{p}[x] \quad\left(\mathbb{F}_{p}:=\mathbb{Z} / p \mathbb{Z}\right.$ is a finite field So $I_{p}[x]$ is a Enclidean domain/has a division algorithm). Thus

$$
f(x) \equiv(x-a) q(x)+r(\bmod p),
$$

Where $g(x) \in F_{p}[x]$ has degree $k$, leading coefficient not divisible by $p$, and $r \in \mathbb{F}_{p}$. Now,

$$
\begin{aligned}
f(a) & \equiv(a-a) g(a)+r(\bmod p) \\
& \equiv r(\bmod p) .
\end{aligned}
$$

Since $f(a) \equiv O(\bmod p)$, we have that $r \equiv O$ riodp ).

Thus $f(x) \equiv(x-a) q(x)(\bmod p)$.

Let $b$ be any solution to $f(x) \equiv 0(\bmod p)$. Then $f(b) \equiv(b-a) q(b)(\bmod p)$.\\
If $b \not \equiv a(\bmod p)$, then $b-a \not \equiv 0(\bmod p)$ and so $q(b) \equiv 0(\bmod p)$. By the induction hypothesis, $q(x)$ has at most $k$ incongruent solutions modulo $p$. Thus $f(x)$ has at most $k+1$ incongruent solutions modulo $p$, as required.

Proposition 5.8: Let $p$ be a prime and let $d \in \mathbb{Z}$ with $d>0$ and $d \mid p-7$. Then the congruence

$$
x^{d}-1 \equiv O(\bmod p)
$$

has exactly $d$ incongruent solutions modulo $p$. Proof: Since $d / p-1$, there exists $e \in \mathbb{Z}$ such that $p-1=d e$. Then

$$
x^{p-1}-1=x^{d e}-1
$$

$$
=\left(x^{d}-1\right)\left(x^{d(e-1)}+x^{d(e-2)}+\ldots+x^{d}+1\right)
$$

The congruence $x^{p-1}-1 \equiv 0(\bmod p)$ has exactly $p-1$ incongruent solutions modulo $p$ (namely $x \equiv 1,2, \ldots, p-1$ ) by Fermat's Little Theorem. Any such solution must be ce solution to either


\begin{equation*}
x^{d}-1 \equiv 0(\bmod p) \tag{1}
\end{equation*}


or


\begin{equation*}
x^{d(e-1)}+x^{d(e-2)}+\cdots+x^{d}+1 \equiv 0(\bmod (p) \tag{2}
\end{equation*}


(since $\mathbb{F}_{p}$ is a finite field, in particular an integral domain). The congruence (2) has at most $d(e-1)=p-1-d$ incongruent solutions modulo $p$ by Theorem 5.7. Thus the congruence (1) has at least $p-1-(p-1-d)=d$ incongruent solutions modulo $p$. This fact combined with Theorem 5.7 implies that the congruence (1) has exactly $d$ incongruent\\
solutions modulo $p$, as desired.

Exercise: Assuming that 17 has a primitive root determine how many incong ment primitive roots there are modulo 17 .

\begin{itemize}
  \item Do the same for modulus 18.
\end{itemize}

Exercise:

\begin{itemize}
  \item Prove that 3 is a primitive root modulo 43
  \item Use the previons part to find all\\
incongrnent integers having order 14 modulo 43.
\end{itemize}

\end{document}